% SUBSECTION 2.5 "Specification and Measurement of Key Constructs" -- DRAFT for advisor, then user (PDF). DASH-FREE.
% Markers %%Pn%%. ALL numbers BIBLE-VERBATIM (grep-confirmed vs docs/Thesis/_tables_from_bible.tex this session):
%   PRisk 0.0001 (L327); US_EPU 0.0124 (L384); GEPU 0.0181 (L443); competition->UncPre 0.0304/0.0302, ->UncRes 0.0008/0.0023,
%   N 12,728/18,492 (L502,L530); in-table UncPre->UncRes 0.0111/0.0230 (L523); CashRatio->Scrutiny 0.7530/0.8519 (L558).
% Cites VERIFIED + in bib: hassan2020 (use 2020, not 2019), baker2016, davis2016 (provisional, fold as-is), hoberg2016 (= our z_log_TotalSimilarity).
%   hoberg2010 + word 'fluidity' DROPPED (our var is total similarity = hoberg2016, NOT fluidity) -> cite hoberg2016 only.
% BOUNDARY: 2.5 reports VALIDITY results (C3,C5) + the 0.7530 validity (C4 measure check). The scrutiny REJECTION VERDICT
%   + the nulls (-0.0000/-0.0056) + the 0.0408 magnitude are 4.1's -> NOT here (P4 is a pure forward-pointer; co-movement QUALITATIVE).
% Phase C self-scan: 'consistent with' kept VERBATIM (C5); F1 NO 'persistent' for competition (time-varying TNIC) -> locus rests on
%   DISCLOSABLE-IN-ADVANCE, time-variation only does the FE-absorption job (not fused); P4.2 'behaves as intended/tracks cash' NOT bald 'valid',
%   'rises in the quarter before' NOT 'around'; RULE-COH-8 (volume != UncQue); no '---'/'--'; no overclaim.

%%P1%%
Before the residual can carry the weight the empirical strategy places on it, it must clear two demands. It must be convergent, moving with established measures of uncertainty, and it must be discriminant, distinct from observable channels that could pass for the signal without being it. We test both in this section. We also flag, and later test, one further possibility: that the pre-announcement rise reflects not the disclosure bind but the volume of analyst attention to cash, a plausible alternative driver we develop below and put to the data as a side analysis in Section~4.1. Throughout, validity is treated as an empirical question answered here, not an assumption.

%%P2%%
The residual is consistent with established measures of uncertainty, though the evidence is supportive rather than decisive. It is positively associated with firm-level political risk, the share of a firm's earnings call devoted to political risk of \citet{hassan2020}, with a coefficient of $0.0001$ that is statistically significant but economically negligible against the residual's own variation; and with economic policy uncertainty, both the United States newspaper-based index of \citet{baker2016} ($0.0124$) and the global, GDP-weighted index of \citet{davis2016} ($0.0181$). Three features keep this leg modest. The tests are one-tailed; the political-risk association, though measured at the firm-quarter level that matches the residual most closely, is economically trivial; and the policy-uncertainty indices are macro series observed monthly, so every firm shares a single value in any period and the association is identified only off aggregate co-movement within a year. We therefore read the residual as consistent with these measures, not as established by them.

%%P3%%
The discriminant evidence is cleaner, and it leads the validity case. Product-market competition, measured by the text-based total product-market similarity of \citet{hoberg2016}, loads on the scripted presentation but not on the call-varying residual: it predicts presentation uncertainty ($0.0304$ and $0.0302$ across specifications, on $18{,}492$ firm-quarters) while its association with the residual is indistinguishable from zero ($0.0008$ and $0.0023$, both insignificant, on the smaller residual sample of $12{,}728$). The reading is a locating one. Competition is a known, disclosable condition of the business that a firm can address in its prepared remarks, so it surfaces there rather than in the unscripted residual the design reserves for the deal-specific signal. The near-zero is a genuine result and not an artifact of construction: competition is not among the controls the residual is netted against, and in this very sample the scripted presentation does carry into the residual ($0.0111$ and $0.0230$), so orthogonality is not mechanical here; moreover, because competition varies within a firm over time, the firm fixed effects do not absorb it, the firm-fixed-effects presentation association remains estimable and strong ($0.0302$), and the matching null in the residual is therefore informative. Presentation uncertainty is a related channel, and we claim no more than that competition does not surface in the residual.

%%P4%%
One channel deserves separate treatment, because it is the most immediate rival to the anticipatory reading. We construct $\mathrm{CashScrutiny}$, the share of analyst question-and-answer turns devoted to cash and liquidity, and an indicator $\mathrm{HighCashScrutiny}$ for calls above the median, which is zero, since most calls draw no cash scrutiny at all. The measure behaves as intended: a firm's cash position predicts how much analysts ask about cash ($0.7530$ and $0.8519$), so it tracks what it is built to track. It is also a plausible alternative driver of the run-up, because cash scrutiny itself rises in the quarter before a cash acquisition, so it co-moves with both the cash position and the event and could in principle generate the pre-announcement rise mechanically. We therefore treat it as an identification threat and test it directly, as a side analysis whose results are reported in Section~4.1. The measure is the volume of cash questions, not their uncertainty wording: the latter, $\mathrm{UncQue}$, is already netted out inside the residual by the decomposition of Section~2.3, so this test addresses a distinct path that the residual does not already absorb.

%%P5%%
The remaining constructs used in the designs are defined here. $\mathrm{CashRatio}$ is the firm's cash-to-assets ratio. $\mathrm{PreAnnounceQtr}$, defined in Section~2.2 and used in Section~2.4, marks the quarter before a firm's first cash acquisition. $\mathrm{CashScrutiny}$ and $\mathrm{HighCashScrutiny}$ are as defined above. The remaining control variables, the firm-financial set of Section~2.4 together with the linguistic and performance controls of the decomposition, are catalogued in the Appendix.
%%END%%
